\section{Related Work}
\textbf{Skill threats and task context.} Our earlier literature review organized the field around threats, benchmarks and defences. We retain that structure but focus it on recovery. Empirical malicious-Skill research and Skill-Inject motivate treating package instructions and helper code as untrusted inputs~\cite{liu,skillinject}. InjecAgent studies indirect injections through tool outputs~\cite{injecagent}. AgentDojo evaluates harmful actions together with legitimate task completion~\cite{agentdojo}. These works motivate both our attack cases and benign controls. Our generated cases do not reproduce their datasets or estimate real attack prevalence.

\textbf{Permissions and information flow.} SkillScope links instruction and code actions, evaluates them under task contexts, and constrains excess privilege~\cite{skillscope}. Progent already creates symbolic task policies and checks whether policy updates narrow or expand authority~\cite{progent}. CaMeL separates trusted control from untrusted data and uses information-flow restrictions~\cite{camel}. These contributions rule out presenting a new task-policy prompt, an action graph or a provenance label as the novelty of this paper. We fix permissions for every method and study what happens after the original route is blocked. Our implementation is a separate finite prototype, not a reproduced full implementation of these systems.

Information-flow control has a long history~\cite{ifc}. LIFTY is especially close: it localizes unsafe accesses and synthesizes policy-enforcing repairs from declarative security specifications~\cite{lifty}. Its guarantees are stronger than our finite testing. Our distinct experimental focus is continuation after partial execution, including a confirmed send, an unknown acknowledgement and an irreversible earlier leak. This is a narrower difference in problem setting, not a claim that automatic security repair is new. The paired-secret test in our verifier checks only a declared finite domain; it is not an unbounded noninterference proof.

\textbf{Construction and feedback.} SKETCH established counterexample-guided synthesis of finite programs~\cite{sketch}. Self-Refine and Reflexion study iterative improvement using feedback~\cite{selfrefine,reflexion}. Our repair and refinement controls are inspired by that general idea but are not exact reproductions of those systems: all use our restricted action grammar and common executable verifier. The proposed constructor adds concrete state transformations before validation. Its ablations test whether those transformations help, while the source-free compiler tests whether the trusted goal already makes them unnecessary.

\textbf{Stateful enforcement and recovery.} Shielding places an external check between an agent and unsafe actions~\cite{shield}. More recent work also joins generated recovery plans with verification, as in PASE~\cite{pase}. MemTX tracks provenance and typed cascading repair for agent memory, explicitly recording leaked irreversible effects~\cite{memtx}. These studies further limit novelty claims based on a generic ``generate, verify, repair'' loop. Here the proposed unit is a residual contract over task predicates, attacker-made dependencies and already executed effects. The evidence must show a benefit for that unit under equal tools, not merely a successful demonstration. Our limited adaptive search and finite fault model leave substantial work before any practical Skill deployment claim.
